\documentclass[12pt]{article}
\usepackage[T2A]{fontenc}
\usepackage[utf8]{inputenc}
\usepackage[russian]{babel}
\usepackage{latexsym,amssymb,amsthm}
\usepackage{amsmath}
\RequirePackage[pdftex]{graphicx}
%%\usepackage{array}

\usepackage{epsf,epic,eepic}
\usepackage{xcolor,bxeepic}
\input ris.tex
\input cells.sty
% \addrisoption{\risleftfalse}

\font\shmt=chess10

\DeclareRobustCommand{\No}{\ifmmode{\nfss@text{\textnumero}}\else\textnumero\fi}

\oddsidemargin=0pt
\textwidth=159mm
\headheight=0pt
\headsep=0pt
\topmargin=0pt
\textheight=240mm

%%
%% Три точки для обозначения делимости
%%
\catcode`\@=11
\def\kratno{\mathrel{\smash{\lower.5ex\hbox{$\,\vdots\,$}}}}
% Это прямые три точки для обозначения делимости
\def\nekratno{\mathrel{\mathpalette\c@ncel\kratno}}
\def\c@ncel#1#2{\m@th\ooalign{$\hfil#1\mkern1mu/\hfil$\crcr$#1#2$}}
% Это перечеркнутые три точки для обозначения неделимости
\def\divis{\smash{\lower.1ex\hbox{\,\hbox{.}\kern-.22em\lower-.6ex\hbox{.}
\kern-.52em\lower-1.2ex\hbox{.}\,}} }
% А это наклонные...

\hyphenation{чис-ло чис-ла чис-лу чис-лом чис-ле чис-лам чис-ла-ми чис-лах
че-ты-рех-уголь-ник че-ты-рех-уголь-ни-ка пря-мо-уголь-ник тре-уголь-ни-ка
тре-уголь-ни-ков вне-впи-сан-ной вне-впи-сан-ных}

\sloppy
\pagestyle{empty}

\def\q#1.{{\bf #1. }}
\def\dotline{\smallskip\hbox to \hsize{\dotfill}\medskip}

\begin{document}

\centerline{\sc САНКТ-ПЕТЕРБУРГСКАЯ}
\centerline{\sc ОЛИМПИАДА ШКОЛЬНИКОВ 2025 года ПО МАТЕМАТИКЕ}
\centerline{\sc II тур. \ 6 класс.}
\medskip
\hrule
\bigskip

\q1. Ковбой, выйдя из салуна, движется вдоль прямой с постоянной скоростью. Одну минуту он идет вперед, затем две минуты назад, потом три минуты вперед и т.\,д. В некотором месте на прямой стоит фонарный столб.
Между первой и второй встречей ковбоя со столбом прошла одна минута. Сколько времени прошло между 2025-й и 2026-й встречей?

\q2. В очереди на медосмотр стоят 100 пациентов со справками. Раз в~минуту врач подходит к какому-то из пациентов и говорит ему: <<У~вас мало справок!>> От этого пациент огорчается и идёт в~конец очереди, попутно раздавая по одной справке всем, кто стоял позади него. Если же у~пациента не хватает справок, чтобы это сделать, медосмотр отменяется для всех. Сегодня врач сумел отменить медосмотр за 19 минут. Докажите, что он мог бы отменить медосмотр и~за одну минуту. (Врач всегда знает, сколько справок имеется у каждого пациента, и сам выбирает, к~кому подходить.)

\q3. Клетчатый квадрат $23\times 23$ разрезали по клеточкам на прямоугольники и пронумеровали их.
Могло ли так оказаться, что площадь первого прямоугольника равна 1 или~2, площадь второго~--- 2~или 3, площадь  третьего --- 3 или 4 и т.\,д?

\setrisbox\vbox{\hsize=60pt\centering\setlength{\unitlength}{12pt}

\begin{picture}(3,3)%(-10,-8)
\put(0,0){\grid(3,3)(1,1)}
\put(1.1,1.1){{\shmt k}}
\put(1.9,1.9){\vector(1,1){.6}}\put(1.1,1.9){\vector(-1,1){.6}}\put(1.1,1.1){\vector(-1,-1){.6}}\put(1.9,1.1){\vector(1,-1){.6}}
\put(1.9,1.5){\vector(1,0){.6}}\put(1.5,1.9){\vector(0,1){.6}}\put(1.1,1.5){\vector(-1,0){.6}}\put(1.5,1.1){\vector(0,-1){.6}}
\end{picture}

\vskip2mm

\begin{picture}(5,5)%(-10,-8)
\put(0,0){\grid(5,5)(1,1)}
\put(2.1,2.1){{\shmt n}}
\put(2.9,2.7){\vector(2,1){1.7}}\put(2.7,2.9){\vector(1,2){.9}}
\put(2.3,2.9){\vector(-1,2){.9}}\put(2.1,2.7){\vector(-2,1){1.6}}
\put(2.9,2.3){\vector(2,-1){1.6}}\put(2.7,2.1){\vector(1,-2){.9}}
\put(2.3,2.1){\vector(-1,-2){.9}}\put(2.1,2.3){\vector(-2,-1){1.6}}
\end{picture}
}

\ris[\rismove(0pt,-2mm)]
\q4. Шахматные фигуры <<конь>> и <<король>> умеют делать ходы восьми разных видов (см. рисунки). Сергей изобрел новую фигуру <<кузнечик>>, которая умеет делать ходы $k$ разных видов. Поставив кузнечика на одну из клеток доски $100\times 100$, Сергей последовательно сделал им $10\,000$ ходов, обойдя кузнечиком все клетки доски по одному разу и в результате вернувшись в~исходную клетку. При каком наименьшем $k$ такое могло произойти?
\endris


\dotline

\centerline{Олимпиада 2025 года. II тур. 6 класс. Выводная аудитория.}
\smallskip

\q5. Шестизначное число поделили с остатком на все натуральные числа от 1000 до 2000. Докажите, что хотя бы один из остатков оказался нечётным.

\q6. На окружности выделено $100k$ дуг, не имеющих общих концов. Дуги раскрашены в $k$ цветов: 100 дуг покрашены в первый цвет, 100 дуг во второй и т.\,д., 100 дуг --- в $k$-й цвет. При этом каждая точка окружности покрашена не более чем в два цвета, но для каждой пары цветов какая-то точка окружности покрашена именно в эти два цвета. Найдите наибольшее возможное значение~$k$.


\clearpage

\centerline{\sc САНКТ-ПЕТЕРБУРГСКАЯ}
\centerline{\sc ОЛИМПИАДА ШКОЛЬНИКОВ 2025 года ПО МАТЕМАТИКЕ}
\centerline{\sc II тур. \ 7 класс.}
\medskip
\hrule
\medskip

\q1. Верно ли, что в любом 120-значном числе без нулей можно так переставить цифры, что полученное число будет делиться на 13?

\q2. В выпуклом четырехугольнике $ABCD$ равны стороны $AB$, $BC$ и~$CD$. Докажите, что если $\angle C = 2\angle A$, то $\angle B = 2\angle D$.

\q3. На доске написаны числа 101, 102, \dots, 200. За одну операцию можно совершить одно из трех действий:

1) увеличить на 1 наименьшее число (любое одно из наименьших, если их
несколько) и уменьшить на 1 наибольшее (тоже любое);

2) стереть с доски одно наибольшее и одно наименьшее числа (тоже по одному, если
какое-то из них встречается несколько раз);

3) умножить все числа на доске на 3.

\noindent
Можно ли за несколько операций сделать все числа на доске равными (разумеется, не стирая их все)?

\q4. На доске $2025\times 2025$ стоит несколько фигур <<косоглазая полуладья>>. Каждая из них бьет одну линию (столбец или строку), но не ту, на которой стоит, а соседнюю. Ни одна из стоящих фигур не бьёт никакую другую. Какое наибольшее количество фигур может стоять на доске?

\dotline

\centerline{Олимпиада 2025 года. II тур. 7 класс. Выводная аудитория.}
\smallskip

\q5. В марсианской школе учатся красные, синие и зелёные марсиане, причём красных марсиан не меньше 80. Каждый день три разноцветных марсианина назначаются дежурными по школе. Через некоторое время оказалось, что 20\,\% учеников побывали на дежурстве по 2 раза, 30\,\%~--- по 3 раза, а остальные~--- по 5 раз. Всех красных марсиан, побывавших на дежурстве 5 раз, наградили кактусом. Докажите, что в школе осталось не менее 70 марсиан без кактуса, побывавших на дежурстве по 5 раз.

\q6. На доске записаны три числа: $2025^{2025}$, $2026^{2026}$, $2027^{2027}$. Петя и~Вася играют в игру. За ход разрешается заменить наименьшее число~$A$ на доске на два таких натуральных числа $B$ и $C$, что $B$ и $C$ не взаимно просты и $B + C = A$ (если наименьших чисел несколько, можно взять любое из них). После замены все простые числа с доски стираются. Проигрывает тот, кто не может сделать ход, первым ходит Петя. Кто выигрывает при правильной игре?

\q7. В треугольнике $ABC$ на стороне $AB$ выбрана точка $M$, на стороне $BC$ точка $N$, на стороне $AC$ точки $K$ и $L$. При этом $AM=KL=CN$, $BM=AK$, $BN=CL$. Докажите, что сумма периметров треугольников $BMN$ и $BKL$ больше периметра $ABC$.


\clearpage

\centerline{\sc САНКТ-ПЕТЕРБУРГСКАЯ}
\centerline{\sc ОЛИМПИАДА ШКОЛЬНИКОВ 2025 года ПО МАТЕМАТИКЕ}
\centerline{\sc II тур. \ 8 класс.}
\medskip
\hrule
\bigskip

\q1. Площадь выпуклого четырехугольника $ABCD$ равна 2. На его стороне $AB$ нашлась такая точка $E$, что площадь треугольника $CDE$ равна 1. Докажите, что на стороне $CD$ найдётся такая точка~$F$, что площадь треугольника $ABF$ равна 1.

\q2. По кругу стоят несколько (не меньше трёх) ненулевых чисел. Может ли оказаться, что для каждых трёх подряд идущих по часовой стрелке чисел $a, b, c$ выполняется равенство $c=a+\frac 1b$?

\q3. В городе живёт $100$ котов и $m$ псов и действует много клубов для животных. Каждый кот и каждый пёс \emph{предпочитает} ровно 100 клубов. При каком наибольшем $m$ они заведомо могут пойти каждый в один из
предпочитаемых клубов так, чтобы ни в одном из клубов не встретились кот и пёс? 

\q4. Александр, будучи ценителем творчества Леонарда Эйлера (1707--1783), взял шестизначное число и поделил его с остатком на все числа от 1707 до 1783. Докажите, что хотя бы один из остатков делится на 3.

\dotline

\centerline{Олимпиада 2025 года. II тур. 8 класс. Выводная аудитория.}
\smallskip

\q5. Дано простое число $p$. Некоторые натуральные числа покрашены в синий цвет. Лягушка прыгает по целым точкам числовой прямой вправо, длина каждого её прыжка~--- натуральное число. Она хочет попасть из точки $-1$ в точку~$p$. Первые два прыжка могут иметь любую натуральную длину, но после второго прыжка лягушка должна оказаться в~точке с синей координатой. После этого все прыжки, начиная с третьего, должны иметь синюю длину. Докажите, что количество возможных маршрутов лягушки кратно~$p$.

\q6. На стороне $AC$ треугольника $ABC$ нашлись точки $P,Q$ такие, что $AQ=AB$, $CP=CB$. Точка $K$ на стороне $AB$ такова, что $BK=AP$, точка $L$ на стороне $CB$ такова, что $BL=CQ$. Докажите, что периметр треугольника $ABC$ меньше суммы периметров треугольников $BKL$ и~$BPQ$.

\q7. Даны натуральные числа $m$ и $n$. Федя пишет в каждой клетке таблицы $m\times n$ вещественное число. После этого Серёжа, дав Феде одну конфету, может выбрать одну линию (строку или столбец) и попросить Федю умножить на $-1$ все числа в этой линии. Серёжа стремится к тому, чтобы в каждой из $m+n$ линий сумма чисел стала неотрицательной. Федя, зная об этом, составил исходную таблицу так, чтобы получить побольше конфет. Сколько конфет сможет гарантированно получить Федя?

\end{document}
